\documentclass[11pt]{article}
\usepackage[margin=1in]{geometry}
\usepackage{amsmath,amssymb,amsthm,mathtools}
\usepackage[colorlinks=true,linkcolor=blue,citecolor=blue,urlcolor=blue]{hyperref}
\newtheorem{theorem}{Theorem}
\newtheorem{lemma}[theorem]{Lemma}
\newtheorem{proposition}[theorem]{Proposition}
\newtheorem{corollary}[theorem]{Corollary}
\theoremstyle{remark}\newtheorem{remark}[theorem]{Remark}
\newcommand{\C}{\mathbb C}
\newcommand{\R}{\mathbb R}
\newcommand{\E}{\mathbb E}
\newcommand{\ip}[2]{\langle #1,#2\rangle}
\newcommand{\norm}[1]{\lVert #1\rVert}
\newcommand{\KG}{K_G^{\mathbb C}}
\newcommand{\ind}{\mathbf 1}
\DeclareMathOperator{\Rea}{Re}
\title{A lower bound of 1.373 for the complex Grothendieck constant}
\author{}
\date{}
\begin{document}
\maketitle
\begin{abstract}
We prove that the classical complex Grothendieck constant exceeds
\(1.373\). The construction subtracts an identity term and the full
real third-chaos projection from a holomorphic Gaussian first-chaos
projection. An anisotropic multiplication inequality bounds the cubic
projection in every dimension. An exact pointwise formula uses the
orthogonality of the two midpoint functions associated with complex
phases. The remaining scalar parameter is controlled on its entire
range by convex tangent bounds. A finite family of auxiliary rational
step weights supplies different bounds on prescribed parameter intervals
for one fixed operator. Exact polynomial exponential moments and an
explicit finite matrix complete the proof.
\end{abstract}

\section{The constant and the result}

For a nonzero complex \(m\)-by-\(n\) matrix \(A=(a_{ij})\), let
\begin{align*}
 S_{\C}(A)&=\max_{|\alpha_i|=|\beta_j|=1}
 \left|\sum_{i,j}a_{ij}\alpha_i\overline{\beta_j}\right|,\\
 V_{\C}(A)&=\sup_{d\ge1}\max_{\norm{u_i}=\norm{v_j}=1}
 \left|\sum_{i,j}a_{ij}\ip{u_i}{v_j}\right|,
\end{align*}
where all vectors belong to \(\C^d\) and
\(\ip{u}{v}=\sum_r u_r\overline{v_r}\). The classical complex
Grothendieck constant is
\(\KG=\sup_{m,n,A\ne0}V_{\C}(A)/S_{\C}(A)\).

\begin{theorem}\label{thm:main}
There is an explicitly specified finite complex matrix \(M\) such that
\[
 \frac{V_{\C}(M)}{S_{\C}(M)}>1.373.
\]
Consequently, \(\KG>1.373\).
\end{theorem}

The classical lower bound approximately \(1.33807\) is attributed to
Davie in Haagerup~\cite{Haagerup} and Friedland--Lim--Zhang~\cite{FLZ}.
Haagerup proves an upper bound approximately \(1.40491\).
Theorem~\ref{thm:main} reduces the width of this reported interval by
approximately \(52\) percent. Heilman~\cite[Section 1.8]{Heilman}
proposes analysis of a cubic perturbation as a possible improvement in
the complex case and leaves that analysis open. Here the operator is
specified with its full real third-chaos projection, and its scalar bound
covers every measurable complex phase assignment. The proof is independent
of any unpublished lower-bound argument.

\section{The Gaussian operator and its parameters}

Let \(\gamma_d\) have density \(\pi^{-d}e^{-\norm z^2}\) on
\(\C^d\), and use a first-linear inner product on
\(L^2(\gamma_d;\C)\). Let \(P\) be orthogonal projection onto
\(\operatorname{span}_{\C}\{z_1,\ldots,z_d\}\). The real coordinates
\[
w=(\sqrt2\Rea z_1,\sqrt2\operatorname{Im}z_1,\ldots,
\sqrt2\Rea z_d,\sqrt2\operatorname{Im}z_d)
\]
are independent standard Gaussians. For a standard real Gaussian
coordinate, define the normalized Hermite polynomials by
\[
H_r(x)=\frac{(-1)^r}{\sqrt{r!}}e^{x^2/2}
\frac{d^r}{dx^r}e^{-x^2/2}.
\]
The complex span of the products
\(\prod_{\ell=1}^{2d}H_{a_\ell}(w_\ell)\) with
\(\sum_\ell a_\ell=j\) is denoted by \(\mathcal H_j(\C^d)\).
Let \(C_3\) be orthogonal projection onto \(\mathcal H_3(\C^d)\).
In particular, \(C_3\) is the complexification of the full real
third-chaos projection. The projections \(P,C_3\) commute with
unitary changes of the complex coordinates.

Throughout the construction the parameters are
\begin{equation}\label{eq:parameters}
\lambda=\frac{53326}{10^6},\qquad
\beta=\frac{499883}{10^6},\qquad D=\frac{3447}{5000}.
\end{equation}
Set
\begin{equation}\label{eq:T}
T=P-\lambda I-\beta C_3.
\end{equation}
Its scalar norm is
\[
\norm T_{\infty\to1}
=\sup_{\norm f_\infty,\norm g_\infty\le1}|\ip f{Tg}|.
\]
Below we prove the dimension-independent bound
\(\norm T_{\infty\to1}\le D\). Later sections construct a finite
matrix with scalar value at most \(D\) and vector value greater than
\(1.373D\).

\section{An anisotropic multiplication inequality}

Write \(z_1=\sqrt s e^{i\phi}\), choosing \(\phi\) arbitrarily on
the null set \(z_1=0\). For real functions of \(s\), use the notation
\(E_s[F]=\int_0^\infty F(s)e^{-s}\,ds\). Let
\(q_R(s)\ge q_I(s)>0\) be measurable weights, and put
\begin{align}
 \mathcal Q(k)&=\int\left[q_R(s)(\Rea(e^{-i\phi}k))^2
 +q_I(s)(\operatorname{Im}(e^{-i\phi}k))^2\right]d\gamma_d,
 \label{eq:Q}\\
 a(s)&=\frac12\left(q_R(s)^{-1}+q_I(s)^{-1}\right),&
 e(s)&=\frac12\left(q_R(s)^{-1}-q_I(s)^{-1}\right).
 \label{eq:inverseweights}
\end{align}
Define
\begin{equation}\label{eq:profiles}
p_4(s)=s^3/6,\qquad p_5(s)=s(s-2)^2/2
\end{equation}
and the symmetric matrices
\begin{align}
G_2&=\begin{pmatrix}
E_s[a(s)(s-1)^2]&E_s[e(s)s(s-1)]/\sqrt2\\
E_s[e(s)s(s-1)]/\sqrt2&E_s[a(s)s^2/2]
\end{pmatrix},\label{eq:G2}\\
G_3&=\begin{pmatrix}
E_s[a p_5]&E_s[e(s)s^2(s-2)]/\sqrt{12}\\
E_s[e(s)s^2(s-2)]/\sqrt{12}&E_s[a p_4]
\end{pmatrix}.\label{eq:G3}
\end{align}
For a real symmetric matrix, \(\lambda_{\max}\) denotes its largest
eigenvalue.

\begin{lemma}\label{lem:weighted}
If all the following quantities are at most one,
\begin{equation}\label{eq:gramconditions}
\begin{gathered}
\beta E_s[a],\quad\beta E_s[s/q_R],\quad\beta E_s[s/q_I],
\quad\beta\lambda_{\max}(G_2),\\
\beta E_s[p_5/q_R],\quad\beta E_s[p_5/q_I],
\quad\beta\lambda_{\max}(G_3),
\end{gathered}
\end{equation}
then, in every dimension and for every \(k\in L^2(\gamma_d;\C)\),
\begin{equation}\label{eq:weighted}
\beta\norm{C_3k}_2^2\le\mathcal Q(k).
\end{equation}
An infinite right-hand side is allowed.
\end{lemma}
\begin{proof}
Regard complex \(L^2\) as a real Hilbert space with inner product
\(\Rea\ip\cdot\cdot\). The inverse pointwise quadratic form is
\[
\mathcal Q^{-1}(v)
=\int\left[a(s)|v|^2+e(s)\Rea(e^{-2i\phi}v^2)\right]d\gamma_d.
\]
Use real orthonormal Hermite bases in the transverse coordinates and
the exact-degree decomposition
\[
\mathcal H_3(\C^d)=\bigoplus_{j=0}^3
\mathcal H_j(\C)\otimes\mathcal H_{3-j}(\C^{d-1}).
\]
Different transverse basis elements, including those of different
degrees, are orthogonal in both the conjugate and nonconjugate products.
Consequently it suffices to analyze each first-coordinate degree.
In that coordinate, orthonormal bases are
\begin{align*}
 j=0:&\quad 1,\\
 j=1:&\quad z_1,\overline{z_1},\\
 j=2:&\quad z_1^2/\sqrt2,
 s-1,\overline{z_1}^2/\sqrt2,\\
 j=3:&\quad z_1^3/\sqrt6,
 z_1(s-2)/\sqrt2,\overline{z_1}(s-2)/\sqrt2,
 \overline{z_1}^3/\sqrt6.
\end{align*}
Their orthogonality and normalizations follow from angular integration
and \(E_s[s^r]=r!\).

Write each basis function as \(R_m(s)e^{im\phi}\). The term
\(a|v|^2\) is diagonal in this basis. The second term retains exactly
the pairs of angular characters satisfying \(m+n=2\). In degree
zero this leaves only the diagonal entry \(E_s[a]\). In degree one
the self-pair \(m=1\) has real and imaginary eigenvalues
\(E_s[s/q_R]\) and \(E_s[s/q_I]\); the \(m=-1\) entry is their
average. In degree two the pair \((0,2)\) gives \(G_2\) on real
coefficients and the same matrix with negated off-diagonal entries on
imaginary coefficients. These matrices have identical eigenvalues;
the unpaired \(m=-2\) entry is a diagonal entry of \(G_2\).
In degree three the self-pair \(m=1\) gives
\(E_s[p_5/q_R]\) and \(E_s[p_5/q_I]\), while the pair \((-1,3)\)
gives \(G_3\) and its off-diagonal sign reversal. The unpaired
\(m=-3\) entry is a diagonal entry of \(G_3\).

Thus \eqref{eq:gramconditions} implies
\(\mathcal Q^{-1}(v)\le\beta^{-1}\norm v_2^2\) throughout the
range of \(C_3\). Weighted Cauchy--Schwarz gives
\[
|\Rea\ip{k}{v}|^2\le\mathcal Q(k)\mathcal Q^{-1}(v)
\le\beta^{-1}\mathcal Q(k)\norm v_2^2.
\]
Taking the supremum over unit \(v\) in that real Hilbert subspace
proves \eqref{eq:weighted}. For \(d=1\), only the first-coordinate
degree three is present; the same sufficient conditions apply.
\end{proof}

\section{The midpoint phase relation}

For \(\mu>0\), let
\begin{equation}\label{eq:phi}
\Phi_\mu(v)=
\begin{cases}
v^2/(2\mu),& |v|\le2\mu,\\
2|v|-2\mu,&|v|\ge2\mu.
\end{cases}
\end{equation}
Equivalently, \(\Phi_\mu(v)=\max_{|r|\le1}(2vr-2\mu r^2)\).

\begin{lemma}\label{lem:pointwise}
For \(q_R\ge q_I>0\), \(\lambda>0\), and real \(v\),
\begin{align}
&\sup_{\substack{|h|^2+|k|^2=1\\\Rea(h\overline k)=0}}
\left[2v\Rea h-2\lambda|h|^2
+q_R(\Rea k)^2+q_I(\operatorname{Im}k)^2\right]\notag\\
&\hspace{30mm}=\max\{q_R,q_I+\Phi_{\lambda+q_I/2}(v)\}.
\label{eq:pointwise}
\end{align}
\end{lemma}
\begin{proof}
By symmetry take \(v\ge0\). Write
\(h=\rho e^{i\theta}\),
\(k=\pm i\sqrt{1-\rho^2}e^{i\theta}\). At \(\rho=0\), the angle
of \(k\) is unrestricted. With \(x=\cos\theta\), the expression
to maximize is
\[
2v\rho x-2\lambda\rho^2
+(1-\rho^2)\bigl[q_R+(q_I-q_R)x^2\bigr],
\qquad 0\le\rho\le1,
\]
and it suffices to take \(0\le x\le1\). If \(q_R=q_I\), its
maximum occurs at \(x=1\), yielding \eqref{eq:pointwise}.
For \(d_0=q_R-q_I>0\), the maximizing angle at \(\rho<1\) satisfies
\[
x=\min\left\{\frac{v\rho}{d_0(1-\rho^2)},1\right\}.
\]
On the first branch, putting \(w=\rho^2\), the value is
\[
q_R-(2\lambda+q_R)w+\frac{v^2w}{d_0(1-w)}.
\]
This is convex in \(w\), and its branch condition
\(v\sqrt w\le d_0(1-w)\) defines an interval starting at zero.
Its maximum is therefore attained at zero or at the boundary with
\(x=1\). On the second branch the value is
\(q_I+2v\rho-(2\lambda+q_I)\rho^2\). This proves the asserted
upper bound. Its first term is attained at \(h=0,k=1\); its second
term is attained by real \(h\) and purely imaginary \(k\).
The endpoint cases follow directly or by continuity.
\end{proof}

For an auxiliary parameter \(0<b<1/2\) satisfying \(2b^2<\beta\), put
\begin{equation}\label{eq:cv}
c=1-2b^2/\beta>0,\qquad A_0=1-2b>0,
\qquad V(s)=\sqrt s\,(A_0+bs).
\end{equation}

\begin{proposition}\label{prop:scalar}
Suppose the positive weights satisfy \eqref{eq:gramconditions} and
\(E_s[q_R]<\infty\). For \(u\ge0\), define
\begin{equation}\label{eq:Bu}
\mathcal B(u)=\lambda-cu+
E_s\max\{q_R(s),q_I(s)+
\Phi_{\lambda+q_I(s)/2}(\sqrt u\,V(s))\}.
\end{equation}
Then in every dimension,
\begin{equation}\label{eq:scalar_reduction}
\norm T_{\infty\to1}\le\sup_{0\le u\le1}\mathcal B(u).
\end{equation}
\end{proposition}
\begin{proof}
Let \(f,g\) be measurable phase functions, and set
\(h=(f+g)/2\), \(k=(f-g)/2\). Then
\(|h|^2+|k|^2=1\) and \(\Rea(h\overline k)=0\) pointwise.
Self-adjointness gives the exact identity
\begin{align*}
\Rea\ip f{Tg}
={}&\norm{Ph}_2^2-\norm{Pk}_2^2
-\lambda(\norm h_2^2-\norm k_2^2)\\
&-\beta(\norm{C_3h}_2^2-\norm{C_3k}_2^2).
\end{align*}
Let \(t=\norm{Ph}_2\), and use a unitary coordinate change to write
\(Ph=tz_1\); when \(t=0\), choose any first coordinate. Applying
Lemma~\ref{lem:weighted} and discarding \(-\norm{Pk}_2^2\) yields
\[
\Rea\ip f{Tg}\le\lambda+t^2-2\lambda\norm h_2^2
-\beta\norm{C_3h}_2^2+\mathcal Q(k).
\]
The function \(\psi=z_1(s-2)/\sqrt2\) is a unit vector in the range
of \(C_3\), so completing a square gives
\[
-\beta\norm{C_3h}_2^2
\le\frac{2b^2t^2}{\beta}+2\sqrt2bt\Rea\ip h\psi.
\]
Also \(t^2=2t\Rea\ip h{z_1}-t^2\). These formulas reduce the
bound to \(\lambda-ct^2\) plus the integral of
\[
2tV(s)\Rea(e^{-i\phi}h)-2\lambda|h|^2
+q_R(\Rea(e^{-i\phi}k))^2
+q_I(\operatorname{Im}(e^{-i\phi}k))^2.
\]
Lemma~\ref{lem:pointwise} bounds this integral by the one in
\eqref{eq:Bu}. Since \(t\le\norm h_2\le1\), the asserted
real-part bound follows. No pointwise alignment of \(h\) with
\(z_1\) is assumed.

Multiplying one original function by a constant phase gives the
absolute-value bound. Finally, every function in the complex unit
disk is the average of two phase functions: pointwise use
\[
re^{i\theta}=\tfrac12e^{i\theta}(r+i\sqrt{1-r^2})
+\tfrac12e^{i\theta}(r-i\sqrt{1-r^2}),\qquad0\le r\le1.
\]
This decomposition is measurable, with an arbitrary phase at zero.
It extends the estimate to both \(L^\infty\) unit balls, proving
\eqref{eq:scalar_reduction}.
\end{proof}

\begin{corollary}\label{cor:family}
Fix \(\lambda,\beta\), and let \(\mathcal I\) be a finite set of
triples \((b_j,q_{R,j},q_{I,j})\), each satisfying the hypotheses
of Proposition~\ref{prop:scalar}. Let \(\mathcal B_j\) be its function
\eqref{eq:Bu}. For the single operator \eqref{eq:T},
\begin{equation}\label{eq:family_reduction}
\norm T_{\infty\to1}\le
\sup_{0\le u\le1}\min_{j\in\mathcal I}\mathcal B_j(u).
\end{equation}
\end{corollary}
\begin{proof}
For a fixed pair of phase functions, first choose one global phase
rotation making its pairing real and nonnegative. Form its midpoint
functions, perform the coordinate alignment in the preceding proof,
and set \(u=\norm{Ph}_2^2\). These choices are made once. The proof
of Proposition~\ref{prop:scalar} then bounds this same real pairing
by \(\mathcal B_j(u)\) for every \(j\in\mathcal I\). Thus it is
bounded by their minimum at this same \(u\). Taking the supremum
over phase functions and using the measurable phase decomposition
from that proposition proves \eqref{eq:family_reduction}.
\end{proof}

The auxiliary triple may therefore be selected after \(u\) is known.
Only a choice among valid inequalities is made; the parameters
\(\lambda,\beta\) and the operator \(T\) are identical in every
inequality. Below one triple is fixed on each entire parameter
interval, so convexity can be applied to that interval using the
same triple at both endpoints.

\section{The rational family and its Gram conditions}

The index set of the family is
\[
\mathcal I=\{12,29,32,34,35,36,37,38,39,40,41,42,43,44,45,46\}.
\]
Write \(b_j=B_j/10^8\), with integer numerators specified by
\[
\begin{array}{c|r@{\qquad}c|r}
j&B_j&j&B_j\\\hline
12&676510&39&8771813\\
29&4210412&40&9162379\\
32&5558035&41&9540392\\
34&6558419&42&9965338\\
35&7059058&43&10352508\\
36&7516257&44&10727882\\
37&7993148&45&11098212\\
38&8408096&46&11429863
\end{array}
\]
These parameters satisfy \(0<b_j<1/2\) and \(2b_j^2<\beta\).
The accompanying data file \texttt{family.json} gives positive
integers \(R_{j,i},I_{j,i}\), \(j\in\mathcal I\),
\(0\le i\le12000\). Define
\begin{equation}\label{eq:weights}
(q_{R,j}(s),q_{I,j}(s))=
\begin{cases}
(R_{j,i},I_{j,i})/10^8,&i/500\le s<(i+1)/500,\quad i<12000,\\
(R_{j,12000},I_{j,12000})/10^8,&s\ge24.
\end{cases}
\end{equation}
Every entry satisfies
\(12576853\le I_{j,i}\le R_{j,i}\le881396009\).
Thus the weights and their inverses are bounded. The SHA-256
identifier of the data file is
\begin{center}\small\texttt{d0d7b6795988bf9846f3dffd55007ddaa906db45edeef1b830cba22570c1f9fe}.
\end{center}
This finite integer table specifies all weights. The certificate uses
only these integers and the rational parameters; no optimization or
interpolation procedure is part of their definition or verification.

\begin{lemma}\label{lem:gram_certificate}
Each of the 16 weight pairs \eqref{eq:weights}, with the common
parameter \(\beta\) in \eqref{eq:parameters}, satisfies all seven
conditions \eqref{eq:gramconditions} strictly.
\end{lemma}
\begin{proof}
For \(r\le3\), use the exact identity
\begin{equation}\label{eq:moments}
\int_a^\infty s^re^{-s}\,ds
=e^{-a}\sum_{k=0}^r\frac{r!}{k!}a^k.
\end{equation}
Subtracting at successive endpoints gives the cell moments. The last
tail moment is used in full. All scalar and matrix-diagonal entries
in \eqref{eq:gramconditions} are rational combinations of these
moments. The squared off-diagonal entries in \(G_2,G_3\) have
rational factors \(1/2,1/12\).

The program \texttt{verify\_gram.py} evaluates these moments with
the rational interval arithmetic in \texttt{arithmetic.py}. It encloses
\(e^{-1/500}\) between the consecutive Taylor sums through degrees
61 and 60, respectively, and obtains the endpoint exponentials by
interval multiplication. Every rational conversion, multiplication,
and division rounds outward onto \(10^{-90}\mathbb Z\), with integer
endpoints. All sums and differences are also enclosed outwardly.
There is no numerical quadrature in this calculation.

Across all family members, each of the five scalar quantities in
\eqref{eq:gramconditions} and each diagonal entry of
\(\beta G_2,\beta G_3\) is less than
\[
0.999989989464410785<1.
\]
For every family member the determinant checks give
\begin{align*}
\det(I-\beta G_2)&>0.053824287051790764,\\
\det(I-\beta G_3)&>0.105885670985068696.
\end{align*}
The diagonal entries of each \(I-\beta G_j\) are positive, and its
determinant is positive. Thus these matrices are positive definite,
which proves the two eigenvalue conditions in
\eqref{eq:gramconditions}. Every displayed decimal endpoint denotes
an exact rational number.
\end{proof}

\section{A certified bound on the entire scalar parameter range}

\subsection{Convex tangent bounds}

Fix one family member and suppress its index for the moment. For fixed
\(s\), put \(a_s=2\lambda+q_I(s)\). The function
\(u\mapsto\Phi_{a_s/2}(\sqrt u\,V(s))\) is concave on
\([0,\infty)\): it is linear up to its transition and then a
positive multiple of \(\sqrt u\) plus a constant, with matching
derivatives at the transition. For any \(\tau>0\), its tangent at
\(u=\tau^2\) is
\begin{equation}\label{eq:tangent}
L_{\tau,u}(V)=
\begin{cases}
uV^2/a_s,&\tau V\le a_s,\\
(\tau+u/\tau)V-a_s,&\tau V\ge a_s.
\end{cases}
\end{equation}
This is a global upper bound for the concave function. For fixed
\(u\ge0\), the two branches agree at their transition and together
form a nondecreasing function of \(V\ge0\).

On \([24,\infty)\) we instead use
\begin{equation}\label{eq:tailbound}
\max\{q_R,q_I+\Phi_{a_s/2}(\sqrt u V)\}
\le q_R+q_I+uV^2/a_s.
\end{equation}
Indeed \(\Phi_{a_s/2}(v)\le v^2/a_s\) on both branches; on the
linear branch the difference is \((|v|-a_s)^2/a_s\).
Define
\begin{align}
J_\tau(u)={}&\lambda-cu+
\int_0^{24}e^{-s}\max\{q_R,q_I+L_{\tau,u}(V(s))\}\,ds\notag\\
&+\int_{24}^\infty e^{-s}(q_R+q_I+uV(s)^2/a_s)\,ds.
\label{eq:J}
\end{align}
Then \(\mathcal B(u)\le J_\tau(u)\). Moreover, \(J_\tau\) is
convex in \(u\): each integrand on \([0,24]\) is the maximum of
two affine functions of \(u\), and the remaining terms are affine.
Consequently,
\begin{equation}\label{eq:convex_cover}
\sup_{u\in[l,r]}\mathcal B(u)
\le\max\{J_\tau(l),J_\tau(r)\}.
\end{equation}

Restore the family index, writing \(J_{j,\tau}\) for
\eqref{eq:J} with the data of member \(j\). Use exactly the
64 intervals \([k/64,(k+1)/64]\), with
\begin{equation}\label{eq:tau}
\tau_k=10^{-6}\left\lfloor10^6\sqrt{\frac{2k+1}{128}}\right\rfloor,
\qquad0\le k<64.
\end{equation}
Assign the following single member to the entire \(k\)-th interval:
\begin{equation}\label{eq:assignment}
\sigma(k)=
\begin{cases}
12,&0\le k\le20,\\
29,&21\le k\le29,\\
32,&30\le k\le32,\\
34,&33\le k\le34,\\
k,&35\le k\le45,\\
46,&46\le k\le63.
\end{cases}
\end{equation}
All \(\tau_k\) are positive rationals. A tangent at this rational
point suffices for \eqref{eq:convex_cover}; it need not be taken
at the exact midpoint. By \eqref{eq:convex_cover} and
Corollary~\ref{cor:family},
\begin{equation}\label{eq:family_cover}
\norm T_{\infty\to1}\le
\max_{0\le k<64}\max\left\{
J_{\sigma(k),\tau_k}(k/64),
J_{\sigma(k),\tau_k}((k+1)/64)\right\}.
\end{equation}
Indeed, on any one interval the minimum in
\eqref{eq:family_reduction} is bounded above by its selected
\(\mathcal B_{\sigma(k)}\), and convexity is applied with that
fixed member and that fixed tangent. The 128 endpoint bounds in
\eqref{eq:family_cover} consequently control every \(u\in[0,1]\).
The assignment \eqref{eq:assignment} and the tangent parameters are
part of the finite data; their validity follows from the exact
inequalities below, without any assumption about how they were chosen.

\subsection{Exact upper integration}

Write \(F(s)=V(s)^2=s(A_0+bs)^2\). The functions \(V,F\) increase
on \([0,\infty)\), since \(A_0,b>0\). On a weight cell \([l,r]\),
the values \(q_R,q_I,a_s\) are positive rational constants.

If \(\tau^2F(r)\le a_s^2\), the whole cell uses the quadratic
branch of \eqref{eq:tangent}. If
\(uF(r)\le a_s(q_R-q_I)\), the maximum in \eqref{eq:J} equals
\(q_R\) throughout the cell. If
\(uF(l)\ge a_s(q_R-q_I)\), it equals \(q_I+uF/a_s\) throughout
the cell. Both cases are integrated exactly by \eqref{eq:moments}.

If \(\tau^2F(l)\ge a_s^2\), the whole cell uses the linear
branch. Put \(L=\tau+u/\tau>0\). The maximum is constantly
\(q_R\) when \(L^2F(r)\le(q_R+2\lambda)^2\), and equals
\(LV(s)-2\lambda\) throughout the cell when
\(L^2F(l)\ge(q_R+2\lambda)^2\). All squared comparisons have
nonnegative sides before squaring.

For the latter case choose the positive rational number
\[
r_0=10^{-8}\left\lfloor10^8\sqrt{(l+r)/2}\right\rfloor.
\]
The global inequality \(\sqrt s\le(s+r_0^2)/(2r_0)\), multiplied
by \(A_0+bs>0\), gives the polynomial upper bound
\begin{equation}\label{eq:Vmajorant}
V(s)\le\frac{A_0r_0}{2}
+\left(\frac{A_0}{2r_0}+\frac{br_0}{2}\right)s
+\frac b{2r_0}s^2.
\end{equation}
Its coefficients are rational, and multiplication by \(L>0\)
followed by \eqref{eq:moments} yields the required integral upper
bound.

For a cell crossing a tangent or maximum transition, subdivide its
length \(h=1/500\) into 16 equal subintervals. Throughout the
original cell the weights remain fixed. On each subinterval,
monotonicity in \(V\) bounds the integrand by the right-endpoint
value of \(\max(q_R,q_I+L_{\tau,u}(V))\). If
\(m_0=\int_l^r e^{-s}\,ds\), the exact mass of its subinterval
indexed by \(v\in\{0,\ldots,15\}\) is
\begin{equation}\label{eq:subcell_mass}
 m_0\frac{e^{-vh/16}(1-e^{-h/16})}{1-e^{-h}}.
\end{equation}
Every factor in this expression is outwardly enclosed using
\(e^{-1/500}\) and \(e^{-1/8000}\). Both exponentials are enclosed
between their Taylor sums through degrees 61 and 60. This expression
for the subinterval mass follows by integrating the exponential over
\([l+vh/16,l+(v+1)h/16]\); it does not assume that the exponential
density is constant.

At a rational subinterval endpoint \(s_*\), the branch is decided
by the exact comparison \(\tau^2F(s_*)\le a_s^2\). A quadratic
endpoint value is rational. A linear endpoint value uses a rational
upper bound on \(V(s_*)=\sqrt{F(s_*)}\), obtained by rounding an
integer square root upward to a multiple of \(10^{-24}\).
Monotonicity gives the upper estimate even when a subinterval itself
contains one or more transitions. The tail in \eqref{eq:J} is
integrated in full by \eqref{eq:moments}, since its weights are
constant and \(V^2\) is cubic.

All square-root decisions are exact integer squaring comparisons.
The choices \(\tau_k,r_0\) need only be positive rational numbers;
their downward rounding does not affect the tangent or polynomial
inequalities. The endpoint upper bound for \(V(r)\) is rounded
upward as required. The interval arithmetic for moments is the same
outward arithmetic used in Lemma~\ref{lem:gram_certificate}.

\begin{lemma}\label{lem:certificate}
For the family \eqref{eq:weights} and the fixed assignment
\eqref{eq:assignment}, the 128 endpoint upper integrals in
\eqref{eq:family_cover} are all less than \(D=3447/5000\).
In particular,
\begin{equation}\label{eq:scalar}
\norm T_{\infty\to1}\le\frac{3447}{5000}
\end{equation}
in every dimension.
\end{lemma}
\begin{proof}
The program \texttt{verify\_scalar.py}, using
\texttt{scalar\_kernel.py} and the integer interval arithmetic in
\texttt{arithmetic.py}, performs the classifications and integrations
above. It verifies the data identifier, the complete 64-interval
partition, the validity of each selected member, and the occurrence
of both endpoints of each interval exactly once with its prescribed
member and tangent. Every endpoint integral in
\eqref{eq:family_cover} is bounded above by an explicitly integrated
majorant. Among the 1,536,000 original radial cell evaluations,
7,304 use the 16 subdivisions in \eqref{eq:subcell_mass}; all other
cells use their polynomial or constant upper integrands directly.

The largest upper enclosure is less than
\[
0.689291686947801369<\frac{3447}{5000}.
\]
It occurs for member 41, interval \(k=41\), right endpoint
\(u=21/32\), and \(\tau=100657/125000\). The margin below \(D\)
is greater than \(0.000108313052198631\). These numbers bound the
computed upper majorants; no lower bound on \(\mathcal B_j\) or
the operator norm is inferred from them. Equation
\eqref{eq:family_cover} now proves the bound for every parameter in
\([0,1]\), and Lemma~\ref{lem:gram_certificate} verifies every
required projection inequality. The programs use Python integers
and exact fractions only. No numerical quadrature, numerical
eigenvalue calculation, or numerical maximization is used in the
verification.
\end{proof}

\section{A feasible vector assignment}

For \(Z\) distributed according to \(\gamma_d\), let
\(R=\norm Z\), \(\mu_d=\E R\), and define
\(F(z)=z/\norm z\) off the null set \(z=0\). The vector function
\(F\) has unit norm almost everywhere.

\begin{lemma}\label{lem:vector}
For the operator \eqref{eq:T}, the assignment \(F\) has value
\begin{equation}\label{eq:Nd}
 N_d:=\sum_{j=1}^d\ip{F_j}{TF_j}
 =\frac{\mu_d^2}{d}\left(1-\frac{\beta}{4(d+1)}\right)-\lambda.
\end{equation}
For the parameters \eqref{eq:parameters},
\begin{equation}\label{eq:NdLower}
 N_d\ge\frac{2d}{2d+1}
 \left(1-\frac{\beta}{4(d+1)}\right)-\lambda,
 \qquad N_d\longrightarrow1-\lambda.
\end{equation}
\end{lemma}
\begin{proof}
Write \(Z=RU\), where \(U\) is uniform on the complex unit sphere
and independent of \(R\). Gaussian polar integration gives
\[
 \E R^2=d,\qquad \E R^3=(d+\tfrac12)\mu_d,
 \qquad \E|U_j|^2=\frac1d,
\]
and
\[
 \E|U_j|^4=\frac2{d(d+1)},\qquad
 \E|U_j|^2|U_k|^2=\frac1{d(d+1)}\quad(j\ne k).
\]
For example, the sphere moments follow by writing the independent
exponential random variables \(|Z_j|^2\) as their sum times the
normalized coordinates; the normalized coordinates have uniform
simplex density. Coordinatewise phase rotations show that
\(\ip{F_j}{z_k}=\delta_{jk}\mu_d/d\). Hence
\(\sum_j\norm{PF_j}_2^2=\mu_d^2/d\).

Under coordinatewise phase rotations, \(F_j\) has character \(e_j\).
A degree-three complex Hermite polynomial with this character has
bidegrees \((2,1)\). The possible orthonormal basis elements are
exactly
\[
 B_j=z_j(|z_j|^2-2)/\sqrt2,
 \qquad B_{jk}=z_j(|z_k|^2-1)\quad(k\ne j).
\]
Indeed, its holomorphic and antiholomorphic multiindices satisfy
\(\alpha-\delta=e_j\), \(|\alpha|+|\delta|=3\), so
\(\delta=e_k\), \(\alpha=e_j+e_k\). Gaussian moments give
\begin{align*}
 \ip{F_j}{B_j}
 &=\frac1{\sqrt2}\left(
 \frac{2(d+1/2)\mu_d}{d(d+1)}-\frac{2\mu_d}{d}\right)
 =-\frac{\mu_d}{\sqrt2\,d(d+1)},\\
 \ip{F_j}{B_{jk}}
 &=\frac{(d+1/2)\mu_d}{d(d+1)}-\frac{\mu_d}{d}
 =-\frac{\mu_d}{2d(d+1)}.
\end{align*}
Squaring and summing over all coordinates gives
\[
 \sum_j\norm{C_3F_j}_2^2
 =d\left(\frac{\mu_d^2}{2d^2(d+1)^2}
 +(d-1)\frac{\mu_d^2}{4d^2(d+1)^2}\right)
 =\frac{\mu_d^2}{4d(d+1)}.
\]
Since \(\sum_j\norm{F_j}_2^2=1\), this proves \eqref{eq:Nd}.
Finally, Cauchy--Schwarz gives
\[
 d^2=(\E R^2)^2\le(\E R)(\E R^3)
 =(d+\tfrac12)\mu_d^2,
 \qquad \mu_d^2\le d.
\]
The positive factor \(1-\beta/[4(d+1)]\) yields
\eqref{eq:NdLower} and the stated limit.
\end{proof}

\section{An explicit finite matrix}

We now fix
\begin{equation}\label{eq:finiteparameters}
 d=10^9,\qquad L=10^9,\qquad \Delta=2^{-80}.
\end{equation}
In each of the \(2d\) standard real coordinates, partition
\([-L,L]\) into consecutive intervals of length \(\Delta\), and
add the two exterior intervals. Assign endpoints consistently; their
Gaussian measures are zero. Let \(\mathcal P\) be the finite product
partition, whose size is \((2L2^{80}+2)^{2d}\).

For two standard real vectors \(w,w'\in\R^{2d}\), write
\(a=w\cdot w'\), \(r=\norm w^2\), \(s=\norm{w'}^2\). The real
degree-three projection kernel is
\begin{equation}\label{eq:kernel}
 H_3(w,w')=\frac{a^3-3a(r+s)+3(2d+2)a}{6}.
\end{equation}
To verify the formula, sum the products of degree-three normalized
Hermite polynomials, separating the multiindices of types \((3)\),
\((2,1)\), and \((1,1,1)\). Their contributions are respectively
\[
 \sum_i\frac{(w_i^3-3w_i)(w_i'^3-3w_i')}{6},\qquad
 \sum_{i\ne j}\frac{(w_i^2-1)(w_i'^2-1)w_jw_j'}2,
 \qquad\sum_{i<j<k}w_iw_jw_kw_i'w_j'w_k'.
\]
Expanding these expressions gives \eqref{eq:kernel}.

For cells \(A,B\in\mathcal P\), define the complex matrix
\begin{equation}\label{eq:matrix}
 M_{AB}=\int_{A\times B}
 \left(\sum_{j=1}^d\overline{z_j}v_j-\beta H_3(w,w')\right)
 \,d\gamma_d(z)d\gamma_d(v)
 -\lambda\gamma_d(A)\ind_{A=B}.
\end{equation}
Here \(w,w'\) are the standardized real coordinates corresponding to
\(z,v\). All entries are explicit finite sums of products of
one-dimensional Gaussian interval moments, with the displayed diagonal
mass term. In particular, \eqref{eq:matrix} specifies a finite matrix
without a limiting definition.

\begin{lemma}\label{lem:finite}
The matrix \eqref{eq:matrix} satisfies
\begin{equation}\label{eq:finitebounds}
 S_{\C}(M)\le D,\qquad
 V_{\C}(M)\ge Q_d:=\frac{2d}{2d+1}
 \left(1-\frac{\beta}{4(d+1)}\right)-\lambda-2\cdot10^{-12}.
\end{equation}
\end{lemma}
\begin{proof}
The coefficient kernel in \eqref{eq:matrix} has the conjugation
appropriate to a first-linear inner product. Directly,
\[
 \ip{f}{Pg}=\iint\left(\sum_j\overline{z_j}v_j\right)
 f(z)\overline{g(v)}\,d\gamma_d(z)d\gamma_d(v).
\]
The kernel \(H_3\) is real symmetric, and the identity term on step
functions is \(\sum_A\gamma_d(A)f_A\overline{g_A}\).
Thus every pair of cellwise phase assignments satisfies
\[
 \sum_{A,B}M_{AB}\alpha_A\overline{\beta_B}=\ip{f}{Tg}.
\]
Corollary~\ref{cor:family} and Lemma~\ref{lem:certificate} prove
the scalar bound in \eqref{eq:finitebounds} for all phases.

For each bounded cell, let \(z_A\) be its complex midpoint and put
\(F_{\mathcal P}(A)=z_A/\norm{z_A}\); for every exterior cell put
\(F_{\mathcal P}(A)=0\). A midpoint is never zero, since every real
coordinate midpoint is an odd multiple of \(2^{-81}\).
On a bounded cell,
\(\norm{z-z_A}\le\sqrt d\,\Delta/2\). For \(\norm z\ge1\),
the normalization inequality gives
\[
 \norm{F(z)-F_{\mathcal P}(z)}
 \le2\norm{z-z_A}/\norm z\le\sqrt d\,\Delta.
\]
On the remaining bounded cells the difference is at most two, while
on exterior cells its norm is one almost everywhere. Markov's
inequality gives
\[
 \Pr(\norm Z<1)\le e\E e^{-\norm Z^2}=e2^{-d}.
\]
The standard real Gaussian fourth moment is three, so a union bound
gives exterior probability at most \(6d/L^4\). Therefore
\begin{equation}\label{eq:approx}
 \norm{F-F_{\mathcal P}}_2^2
 \le d2^{-160}+4e2^{-d}+6d/L^4<10^{-24}.
\end{equation}
For the last comparison, the first term is below \(10^{-30}\)
because \(d10^{30}<2^{160}\). The second is below \(10^{-30}\)
using \(e<3\), \(d\ge200\), and \(12\cdot10^{30}<2^{200}\).
The third is \(6\cdot10^{-27}\).

The ranges of \(P\) and \(C_3\) are orthogonal. The eigenvalues of
\(T\) are \(1-\lambda\), \(-\lambda-\beta\), and \(-\lambda\),
on their respective orthogonal subspaces. Since
\(\lambda+\beta=553209/10^6<1\),
\[
 \norm{T}_{2\to2}=\max(1-\lambda,\lambda+\beta)<1.
\]
The same norm holds for componentwise action on vector-valued
\(L^2\). Since \(\norm{F}_2=1\), \(\norm{F_{\mathcal P}}_2\le1\),
\eqref{eq:approx} implies
\[
 \left|\sum_j\ip{F_j}{TF_j}
 -\sum_j\ip{(F_{\mathcal P})_j}{T(F_{\mathcal P})_j}\right|
 <2\cdot10^{-12}.
\]
Finally assign to the row and column cells the unit vectors
\[
 u_A=\left(F_{\mathcal P}(A),
 \sqrt{1-\norm{F_{\mathcal P}(A)}^2},0\right),\qquad
 v_B=\left(F_{\mathcal P}(B),0,
 \sqrt{1-\norm{F_{\mathcal P}(B)}^2}\right).
\]
Their cross inner products are those of \(F_{\mathcal P}\).
The real part of their matrix value is therefore at least \(Q_d\),
by Lemma~\ref{lem:vector}. This proves the vector bound.
\end{proof}

\begin{proof}[Proof of Theorem~\ref{thm:main}]
The scalar constant is \(D=3447/5000\). For the explicit finite
parameters \eqref{eq:finiteparameters}, the rational numerator lower
bound in \eqref{eq:finitebounds} is
\[
Q_d=\frac{315558000264346749990112333333}
{333333333833333333500000000000}.
\]
An exact rational comparison gives
\[
Q_d-\frac{1373}{1000}D
=\frac{42599791073649832354633333}
{333333333833333333500000000000}>0.
\]
In particular \(Q_d>0\), so \(M\ne0\). Lemma~\ref{lem:finite}
yields
\(V_{\C}(M)/S_{\C}(M)\ge Q_d/D>1373/1000=1.373\).
The corresponding limiting comparison is
\[
\frac{1-\lambda}{D}=\frac{52593}{38300}>1.373.
\]
\end{proof}

\begin{thebibliography}{9}
\bibitem{Haagerup}
U.~Haagerup, \emph{A new upper bound for the complex Grothendieck
constant}, Israel Journal of Mathematics \textbf{60} (1987), no.~2,
199--224. \href{https://doi.org/10.1007/BF02790792}{doi:10.1007/BF02790792}.

\bibitem{FLZ}
S.~Friedland, L.-H.~Lim, and J.~Zhang,
\emph{An elementary and unified proof of Grothendieck's inequality},
L'Enseignement Math\'ematique \textbf{64} (2018), no.~3/4, 327--351.
\href{https://doi.org/10.4171/LEM/64-3/4-6}{doi:10.4171/LEM/64-3/4-6}.

\bibitem{Heilman}
S.~Heilman, \emph{A lower bound for Grothendieck's constant},
\href{https://arxiv.org/abs/2603.22616v1}{arXiv:2603.22616v1} (2026).
\end{thebibliography}
\end{document}
